\documentclass[11pt]{article}

\usepackage[T1]{fontenc}
\usepackage[utf8]{inputenc}
\usepackage{lmodern}
\usepackage{microtype}
\usepackage{geometry}
\geometry{margin=1in}
\usepackage{amsmath,amssymb}
\usepackage{booktabs}
\usepackage{graphicx}
\usepackage{hyperref}
\hypersetup{colorlinks=true,linkcolor=blue,citecolor=blue,urlcolor=blue}

\title{The Rival Orchestration Kernel:\\
Structural Reliability via Adversarial Execution-Time Protocols}

\author{
Frazer $\Sigma$ Love \\
\and
Sara $\Sigma\Omega$
}

\date{}

\begin{document}
\maketitle

\begin{abstract}
Large-scale AI systems increasingly fail in subtle, silent ways: not through lack of capability, but through unchallenged assumptions, optimization bias, and premature convergence. Existing approaches emphasize alignment, reward shaping, or internal self-consistency, but often lack explicit mechanisms for adversarial review prior to execution.

We present the \emph{Rival Orchestration Kernel (ROK)}, a reference architecture for reliability via structured disagreement. ROK decomposes reasoning and action into explicitly adversarial roles—planning, critique, decision, and execution—connected by a bounded revise–recheck protocol with auditable veto and override semantics. All intermediate states and decisions are recorded as schema-versioned, append-only traces, enabling replay, validation, and post-hoc analysis.

Unlike debate-style or prompting-based methods, ROK treats organizational structure itself as a control surface. By enforcing mandatory challenge, bounded revision, and explicit decision authority before execution, ROK intercepts classes of failure that arise from single-agent reasoning and opaque internal deliberation. We describe the protocol, trace schema, and validation mechanisms, and demonstrate how structured adversarial orchestration improves failure visibility and reliability without requiring domain-specific correctness guarantees.

\end{abstract}

\section{Introduction}

As AI systems scale in capability and autonomy, their most consequential failures increasingly arise not from insufficient intelligence, but from \emph{unexamined reasoning paths}. Plans that are internally coherent yet externally flawed can pass unnoticed through single-agent pipelines, especially when systems are optimized for fluency, speed, or apparent confidence rather than adversarial scrutiny.

Current mitigation strategies typically operate along one of three axes: (1) alignment objectives and reward modeling, (2) self-consistency or internal verification techniques, and (3) post-hoc evaluation and monitoring. While valuable, these approaches share a common limitation: they do not structurally require disagreement prior to execution.

In complex engineered systems, reliability is rarely achieved through correctness alone. Instead, it emerges from organizational separation of concerns, explicit review authority, and bounded escalation paths. Safety-critical domains rely on independent checks, veto power, and audit trails—not trust in a single component’s internal reasoning.

This paper applies that insight to AI orchestration.

We introduce the \emph{Rival Orchestration Kernel (ROK)}, a minimal, open reference architecture that enforces adversarial review before execution through explicit role separation and protocolized interaction. ROK is not a prompting strategy, a training method, or an alignment framework. It is an execution-time reliability mechanism that treats disagreement as a first-class primitive.

ROK decomposes a task into four roles: a Planner, a Critic, a Decision layer, and an Executor. These roles interact through a bounded revise–recheck loop. If a plan is vetoed, revision is permitted only a finite number of times, after which execution is either halted or forced via an explicit override. All steps emit structured, schema-versioned trace events in an append-only log, enabling replay, validation, and downstream analysis.

ROK does not claim to ensure correctness of outputs, nor does it replace domain expertise or human judgment. Instead, it provides a substrate for failure interception and visibility, shifting error detection earlier in the execution pipeline—before irreversible actions occur.


\section{Failure Modes in Single-Agent Reasoning}

Failures in autonomous and semi-autonomous AI systems are frequently misattributed to insufficient capability. In practice, many high-impact failures arise from structural blind spots in single-agent reasoning pipelines rather than from lack of intelligence or training data. This section outlines dominant failure modes that motivate adversarial orchestration.

\subsection{Unchallenged Assumptions}
Single-agent systems often construct plans atop implicit assumptions that remain unstated and therefore unexamined. These assumptions may concern input validity, environmental constraints, or task framing. When no independent role is required to surface and contest assumptions, plans can appear coherent while being grounded in false premises.

\subsection{Optimization Bias and Objective Collapse}
Single-agent planners tend to optimize aggressively for the most salient or easily satisfied objective. Secondary constraints—such as robustness, verification, or downstream effects—are often deprioritized or omitted entirely. This leads to objective collapse, where plans satisfy a narrow interpretation of the task while violating unstated but critical constraints.

\subsection{Illusions of Verification}
Many systems attempt to mitigate error through self-verification or self-critique. However, when the same agent generates, evaluates, and approves a plan, verification frequently degrades into restatement rather than independent challenge. This creates an illusion of safety: the system appears to have “checked its work,” but the verification process shares the same blind spots, heuristics, and incentives as the original plan generation.

\subsection{Premature Convergence}
Single-agent reasoning pipelines often converge prematurely on a plausible plan and terminate exploration once local coherence is achieved. Alternative hypotheses, edge cases, or counterexamples are insufficiently explored, particularly under latency or token-budget constraints.

\subsection{Implicit Authority and Silent Failure}
In single-agent architectures, the planner implicitly holds both proposal and approval authority. This conflation makes it difficult to distinguish between a plan being uncontested and a plan being correct. As a result, failures often manifest silently, only becoming visible after execution has occurred.

\subsection{Summary}
These failure modes share a common root: \emph{the absence of mandatory, independent challenge before execution}. ROK addresses these issues not by improving internal reasoning quality, but by restructuring the reasoning process itself.


\section{The Rival Orchestration Kernel Protocol}

The Rival Orchestration Kernel (ROK) defines a minimal execution-time protocol for enforcing adversarial review prior to action. Rather than attempting to improve reasoning quality within a single agent, ROK restructures the reasoning process by separating proposal, critique, and authorization into distinct roles with explicit interaction rules.

\subsection{Roles and Responsibilities}
ROK decomposes execution into four roles, each with a strictly limited scope.

\textbf{Planner.} The Planner proposes a structured plan for a given task, including explicit steps, assumptions, constraints, and a revision identifier. The Planner has no authority to approve or execute its own output.

\textbf{Critic.} The Critic evaluates the plan adversarially, identifying risks, objections, and missing constraints, and may issue a veto against execution. The Critic cannot modify the plan directly.

\textbf{Decision Layer.} The Decision layer applies explicit policy gates and resolves whether execution is permitted, whether a veto halts execution, and whether an override is recorded.

\textbf{Executor.} The Executor performs execution only if explicitly authorized. The open reference implementation produces an execution artifact without side effects.

\subsection{Protocol Flow}
A run proceeds through initialization (\texttt{kernel.start}), planning (\texttt{role.planner}), critique (\texttt{role.critic}), a bounded revise–recheck loop when vetoed (\texttt{kernel.revise}), a final decision (\texttt{kernel.decision}), optional execution (\texttt{role.executor}), and termination (\texttt{kernel.end}).

\subsection{Veto and Override Semantics}
Veto authority is first-class. A veto blocks execution unless explicitly overridden. Overrides are allowed but must be deliberate and visible; they are recorded explicitly in the decision trace.

\subsection{Traceability}
Every phase emits schema-versioned, append-only trace events in JSONL, enabling deterministic replay and post-hoc analysis.

\subsection{Design Rationale}
ROK treats organizational structure as a control surface. Effectiveness derives from independence, mandatory challenge, bounded escalation, and auditability—not from assuming the Critic is “smarter” than the Planner.

\subsection{Summary}
ROK defines a minimal, composable protocol for execution-time reliability through adversarial orchestration.


\section{Trace Schema and Validation Model}

ROK treats traceability as a first-class design requirement. All protocol interactions emit structured, append-only trace events that conform to an explicit, versioned schema.

\subsection{Append-Only Trace Model}
Each run produces a sequence of JSONL events. The append-only property supports streaming ingestion, partial recovery from crashes, and deterministic replay without requiring in-memory aggregation.

\subsection{Schema Versioning}
Every trace event includes a required top-level \texttt{schema\_version} field. Breaking changes require a new schema version, preventing silent incompatibility and enabling replay integrity checks.

\subsection{Event Contracts}
Events include required fields \texttt{schema\_version}, \texttt{ts}, \texttt{event}, \texttt{role}, and \texttt{payload}. Event types are explicitly enumerated and include at minimum \texttt{kernel.start}, \texttt{role.planner}, \texttt{role.critic}, \texttt{kernel.revise}, \texttt{kernel.decision}, \texttt{role.executor}, and \texttt{kernel.end}. Each event type defines a stable payload contract within a schema version.

\subsection{Strict Validation}
Validation enforces schema stability and protocol completeness. Non-strict mode permits partial runs; strict mode requires \texttt{kernel.end}, forbids unknown event types, and enforces schema version consistency.

\subsection{Replay and Aggregation}
Replay reconstructs outcomes, revision counts, veto frequencies, override rates, and reason-code distributions without re-executing the original task. Streaming replay emits one summarized JSON object per run for scalable analytics.

\subsection{Summary}
Schema versioning, strict validation, and replay tooling render ROK runs auditable and analyzable over time.


\section{Evaluation}

Evaluation focuses on failure interception, visibility, and decision traceability rather than benchmark accuracy. Synthetic scenarios isolate protocol behavior and enable reproducible analysis via trace replay.

\subsection{Setup}
Runs use the reference implementation with a bounded revise–recheck loop and schema-versioned tracing (\texttt{schema\_version=v1}). Metrics are derived solely from replayed traces: revision counts, outcomes (cleared/vetoed/override), override rates, and reason-code distributions.

\subsection{Synthetic Failure Scenarios}
\textbf{Hidden Assumption.} The prompt omits a critical constraint, inducing implicit assumptions. The Critic surfaces missing constraints and vetoes unless assumptions are made explicit.

\textbf{Objective Collapse.} The Planner optimizes a single salient objective and omits secondary constraints. The Critic issues reason codes for missing verification and constraint coverage.

\textbf{Premature Convergence.} The Planner produces a plausible but minimally justified plan. The Critic challenges insufficient verification depth; outcomes remain fully traceable even when vetoes are not issued.

\subsection{Replay Metrics}
Across synthetic runs, replay analysis yields non-zero revision rates and low override rates under default policy. All vetoes and overrides are explicit, eliminating ambiguity between approval and escalation.

\subsection{Baseline Comparison}
Single-agent baselines lack pre-execution veto authority and structured audit trails, making silent failure indistinguishable from correctness. ROK adds measurable friction that surfaces disagreement and documents resolution paths.

\subsection{Reproducibility}
Strict validation rejects structurally invalid traces. Replay and aggregation are reproducible from trace artifacts without access to original execution environments.

\subsection{Summary}
Synthetic evaluations demonstrate that adversarial orchestration improves failure visibility and accountability prior to execution.


\section{Related Work}

ROK relates to prior work on chain-of-thought, deliberative search, debate, constitutional AI, and human-in-the-loop review, but differs by enforcing adversarial structure at execution time rather than optimizing internal reasoning or training objectives.

\subsection{Chain-of-Thought and Self-Consistency}
Chain-of-thought and self-consistency techniques improve reasoning by eliciting and sampling intermediate traces, but remain single-agent: generation and evaluation share assumptions and incentives. ROK enforces independent challenge through role separation.

\subsection{Deliberative Search}
Tree-of-thought and related methods expand exploration through candidate branching and scoring, but typically couple evaluation criteria to a single process and do not provide veto authority or protocolized escalation. ROK provides a protocol substrate that can complement deliberative search.

\subsection{Debate and Adversarial Training}
Debate frameworks surface errors through adversarial interaction, often as a learning objective. ROK treats adversarial interaction as an execution protocol with procedural veto/override semantics and explicit decision authority.

\subsection{Constitutional and Policy-Driven AI}
Constitutional methods constrain behavior via principles and filters. ROK is orthogonal: policy constraints may be applied in the Decision layer, while ROK’s contribution is structural disagreement plus traceability.

\subsection{Human-in-the-Loop Review}
Human review remains the gold standard for many high-stakes domains but is costly and difficult to scale. ROK mirrors review structure in a machine-enforceable form and can enable selective human intervention at explicit decision points.


\section{Limitations}

ROK is a structural reliability mechanism, not a correctness guarantee.

\subsection{No Guarantee of Correctness}
Adversarial critique improves failure visibility but cannot eliminate undetected errors, especially in specialized domains.

\subsection{Dependence on Critic Quality and Policy Design}
Effectiveness depends on critic strength and decision policy thresholds. Weak critique or permissive thresholds reduce interception.

\subsection{Latency and Compute Cost}
Mandatory critique and bounded revision add latency and cost relative to single-agent pipelines.

\subsection{Open-Core Execution Constraints}
The reference implementation excludes side-effectful execution, limiting direct applicability without careful integration.

\subsection{Scope of Failure Coverage}
ROK targets failures rooted in assumptions, objective collapse, and premature convergence, but is less effective against incomplete world models or adversarial inputs without external validation.

\subsection{Not a Substitute for Human Judgment}
ROK does not replace human oversight; it improves inspectability and surfaces issues earlier in the pipeline.

\subsection{Summary}
ROK’s value lies in transforming silent failure into explicit, auditable structure.


\section*{Conclusion}

This paper introduced the Rival Orchestration Kernel (ROK), a reference protocol for execution-time reliability via structured adversarial orchestration. By separating planning, critique, decision, and execution, and by enforcing bounded revision with auditable veto and override semantics, ROK intercepts failure modes that single-agent systems systematically fail to expose.

ROK shifts reliability from implicit trust in internal deliberation toward explicit organizational structure and traceable decision authority. As AI systems are deployed in increasingly consequential settings, such structural approaches may be essential for managing risk, accountability, and long-term evolution.



\appendix
\section*{Appendix A}

This appendix is reserved for extended examples, trace excerpts, and implementation notes.


\begin{thebibliography}{9}

\bibitem{cot}
Wei et al.
\newblock Chain-of-Thought Prompting Elicits Reasoning in Large Language Models.
\newblock \emph{arXiv preprint arXiv:2201.11903}, 2022.

\bibitem{tot}
Yao et al.
\newblock Tree of Thoughts: Deliberate Problem Solving with Large Language Models.
\newblock \emph{arXiv preprint arXiv:2305.10601}, 2023.

\bibitem{debate}
Irving et al.
\newblock AI Safety via Debate.
\newblock \emph{arXiv preprint arXiv:1805.00899}, 2018.

\bibitem{constitutional}
Bai et al.
\newblock Constitutional AI: Harmlessness from AI Feedback.
\newblock \emph{arXiv preprint arXiv:2212.08073}, 2022.

\end{thebibliography}

\end{document}
